\section{Future Work}
\label{sec:future-work}

The delivered scope is complete: the four-count DM detector, calibrated scoring,
guarded adaptation, state recovery, streaming interface, local edge simulation,
and labelled development/final evaluations have retained code and evidence.
The following extensions are separate research or deployment work, rather than
unfinished requirements of that deliverable. Earlier proposals and plans are
historical design records; this section defines the remaining optional scope.

\subsection{Broader external validation and temporal stability}

New labelled device traffic, longer benign operating periods, additional attack
intervals, and capture-block confidence intervals would test generalisation
beyond the two labelled references and their small final attack sets. Labels
must remain outside fitting, calibration, and runtime learning.

The separate temporal branch already contains expanding folds, capture/episode
reports, and source-timestamp validation. Its historical findings are retained
in Appendix~\ref{sec:temporal-research}. Integrating those tools with the selected
scoring and both adaptive policies would measure stability across training
boundaries. Any new numerical study using D-Link fitting data should resolve
the known October 8/9 overlap; D-Link supplies no labelled detection claim here.
The completed final runner validates actual chronology of the labelled final
captures, without claiming a merged repository-wide temporal implementation.

\subsection{Adaptation and window duration}

Further drift experiments should test whether conservative updates can improve
legitimate behaviour changes while retaining attack evidence. Rejected candidates
are useful outcomes; weakening guards just to obtain promotions is not a valid
objective. Extra buffer/cadence comparisons, broader baselines, and packet-level
features require a common protocol and compatible evidence across devices.

Thirty-minute static operation merits further study when slower response is
acceptable. The present 10/20/30/60-minute results do not justify replacing the
default across all modes: coarse-window threshold updates increase Samsung false
alerts, and hourly calibration has fewer observations. Temporal false-alert
stability and response-time objectives should guide any future change.

\subsection{Physical deployment and academic submission}

Physical/ARM gateway measurements, a live packet collector, durable event
transport, and downstream response integration extend the supplied-window
Linux/amd64 simulation. They are not measured deployment capabilities of this
release. The existing service and prepared image provide their starting point.

The separate \path{thesis/} workspace supports subsequent dissertation writing.
Completing its chapters, personal metadata, and final PDF/A validation is an
academic submission task; the full project report is the completed technical
and experimental record. Appendix~\ref{sec:thesis-workspace} documents that
workspace and its template attribution without presenting a skeleton as a
finished dissertation.
